% Draf terjemahan Bahasa Indonesia dari chapter2.tex baris 1245--1563.
% Sumber: Methods of Algebra, Volume 2, commit
% 9a5803ff77dd3257484cb177f851a73770a59dd3 (CC BY 4.0).
% stable-unit-id: o014.aljabr2.chapter2.exact-functors-injective-projective-objects
% Segmen ini berhenti tepat sebelum section sumber pada baris 1564.

% segment-id: o014.li2.chapter2.exact-functors-injective-projective-objects.h001
\section{Funktor Eksak, Objek Injektif, dan Projektif}\label{sec:inj-proj}
\hypertarget{o014.aljabr2.chapter2.exact-functors-injective-projective-objects}{ }

% segment-id: o014.li2.chapter2.exact-functors-injective-projective-objects.p001
Untuk suatu funktor $F: \mathcal{A} \to \mathcal{B}$, kita dapat menanyakan apakah
ia mempertahankan $\varinjlim$ atau $\varprojlim$; lihat \cite[\S 2.8]{Li1} atau
pembahasan pada \S\ref{sec:limit-functor}. Bagian ini berfokus pada kasus ketika
$\mathcal{A}$ dan $\mathcal{B}$ merupakan kategori abelian dan $F$ merupakan funktor
aditif. Dengan mengambil kasus khusus $\varprojlim$ (atau $\varinjlim$), yakni
$\Ker$ (atau $\Coker$), untuk setiap morfisme $f:X\to Y$ di $\mathcal{A}$ kita
memperoleh morfisme kanonik $F\Ker(f)\to\Ker F(f)$ dan versi dualnya
$\Coker F(f)\to F\Coker(f)$; keduanya membuat diagram berikut komutatif.

% segment-id: o014.li2.chapter2.exact-functors-injective-projective-objects.g001
\[\begin{tikzcd}[column sep=small]
	F \Ker(f) \arrow[rd, "{F[\Ker(f) \hookrightarrow X]}"'] \arrow[rr] & & \Ker F(f) \arrow[hookrightarrow, ld] \\
	& FX &
\end{tikzcd} \quad \begin{tikzcd}[column sep=small]
	\Coker F(f) \arrow[rr] & & F \Coker(f) \\
	& FY \arrow[twoheadrightarrow, lu] \arrow[ru, "{ F[Y \twoheadrightarrow \Coker(f)] }"'] . &
\end{tikzcd}\]

% segment-id: o014.li2.chapter2.exact-functors-injective-projective-objects.l001
\begin{itemize}
% segment-id: o014.li2.chapter2.exact-functors-injective-projective-objects.i001
	\item Jika $F\Ker(f)\rightiso\Ker F(f)$ berlaku untuk setiap $f$, kita mengatakan
	\emph{$F$ mempertahankan kernel};
% segment-id: o014.li2.chapter2.exact-functors-injective-projective-objects.i002
	\item jika $\Coker F(f)\rightiso F\Coker(f)$ berlaku untuk setiap $f$, kita mengatakan
	\emph{$F$ mempertahankan kokernel}.
\end{itemize}

% segment-id: o014.li2.chapter2.exact-functors-injective-projective-objects.p002
Jika $(X^\bullet,d^\bullet)$ merupakan kompleks dalam kategori abelian, maka
$(FX^\bullet,Fd^\bullet)$ juga demikian; persoalannya adalah keeksakan.

% segment-id: o014.li2.chapter2.exact-functors-injective-projective-objects.pr001
\begin{proposisi}\label{prop:left-right-exact-functor}
Untuk funktor aditif $F:\mathcal{A}\to\mathcal{B}$ antara kategori abelian,
pernyataan berikut ekuivalen:
% segment-id: o014.li2.chapter2.exact-functors-injective-projective-objects.l002
\begin{enumerate}[(L1)]
% segment-id: o014.li2.chapter2.exact-functors-injective-projective-objects.i003
	\item $F$ mempertahankan kernel;
% segment-id: o014.li2.chapter2.exact-functors-injective-projective-objects.i004
	\item jika $0\to X'\xrightarrow{f}X\xrightarrow{g}X''$ eksak di $\mathcal{A}$,
	maka $0\to F(X')\xrightarrow{Ff}F(X)\xrightarrow{Fg}F(X'')$ eksak di $\mathcal{B}$;
% segment-id: o014.li2.chapter2.exact-functors-injective-projective-objects.i005
	\item $F$ mempertahankan semua $\varprojlim$ berhingga.
\end{enumerate}
Secara dual, pernyataan berikut juga ekuivalen:
% segment-id: o014.li2.chapter2.exact-functors-injective-projective-objects.l003
\begin{enumerate}[(R1)]
% segment-id: o014.li2.chapter2.exact-functors-injective-projective-objects.i006
	\item $F$ mempertahankan kokernel;
% segment-id: o014.li2.chapter2.exact-functors-injective-projective-objects.i007
	\item jika $X'\xrightarrow{f}X\xrightarrow{g}X''\to0$ eksak di $\mathcal{A}$,
	maka $F(X')\xrightarrow{Ff}F(X)\xrightarrow{Fg}F(X'')\to0$ eksak di $\mathcal{B}$;
% segment-id: o014.li2.chapter2.exact-functors-injective-projective-objects.i008
	\item $F$ mempertahankan semua $\varinjlim$ berhingga.
\end{enumerate}
\end{proposisi}

% segment-id: o014.li2.chapter2.exact-functors-injective-projective-objects.pf001
\begin{bukti}
Berkat dualitas, cukup membahas (L1)--(L3).

(L1) $\implies$ (L2). Barisan eksak berbentuk
$0\to\bullet\to\bullet\to\bullet$ tepat bersesuaian dengan kernel suatu morfisme;
lihat penjelasan di bagian akhir \S\ref{sec:cohomology}.

(L2) $\implies$ (L3). Semua $\varprojlim$ berhingga dapat dibangun dari produk
berhingga dan penyama. Kita telah mengetahui bahwa funktor aditif
mempertahankan biproduk; karena $F$ mempertahankan $\Ker$, ia mempertahankan
semua penyama (Catatan~\ref{rem:equalizer-Ker}). Jadi, $F$ mempertahankan semua
$\varprojlim$ berhingga.

(L3) $\implies$ (L1) bersifat langsung.
\end{bukti}

% segment-id: o014.li2.chapter2.exact-functors-injective-projective-objects.p003
Untuk funktor aditif $F$ seperti di atas, jika setiap barisan eksak
$(X^\bullet,d^\bullet)$ menghasilkan barisan eksak $(FX^\bullet,Fd^\bullet)$,
kita mengatakan bahwa $F$ mempertahankan barisan eksak. Demikian pula, kita
dapat menanyakan apakah $F$ mempertahankan barisan eksak pendek; ternyata kedua
sifat ini ekuivalen.

% segment-id: o014.li2.chapter2.exact-functors-injective-projective-objects.pr002
\begin{proposisi}\label{prop:exact-functor}
Untuk funktor aditif $F:\mathcal{A}\to\mathcal{B}$ antara kategori abelian,
pernyataan berikut ekuivalen:
% segment-id: o014.li2.chapter2.exact-functors-injective-projective-objects.l004
\begin{enumerate}[(E1)]
% segment-id: o014.li2.chapter2.exact-functors-injective-projective-objects.i009
	\item $F$ mempertahankan barisan eksak pendek;
% segment-id: o014.li2.chapter2.exact-functors-injective-projective-objects.i010
	\item $F$ mempertahankan barisan eksak;
% segment-id: o014.li2.chapter2.exact-functors-injective-projective-objects.i011
	\item $F$ mempertahankan semua $\varprojlim$ berhingga dan $\varinjlim$ berhingga.
\end{enumerate}
\end{proposisi}

% segment-id: o014.li2.chapter2.exact-functors-injective-projective-objects.pf002
\begin{bukti}
(E1) $\implies$ (E2). Diberikan barisan eksak $(X^\bullet,d^\bullet)$ di
$\mathcal{A}$, uraikan ia menjadi barisan eksak pendek
% segment-id: o014.li2.chapter2.exact-functors-injective-projective-objects.g002
\[\begin{tikzcd}
	0 \arrow[r] & \Ker(d^n) \arrow[r, "{\iota^n}"] & X^n \arrow[r, "{\mathbf{d}^n}"] & \Ker(d^{n+1}) \arrow[r] & 0
\end{tikzcd}\]
dengan morfisme $\mathbf{d}^n$ dicirikan oleh
$d^n=\iota^{n+1}\mathbf{d}^n$; indeks atas $n$ boleh dipilih sembarang, tetapi
$X^n$ tidak boleh merupakan suku ujung kanan dari barisan eksak.

Berdasarkan asumsi, barisan
$0\to F(\Ker(d^n))\xrightarrow{F\iota^n}F(X^n)
\xrightarrow{F\mathbf{d}^n}F(\Ker(d^{n+1}))\to0$ tetap eksak, dan
$Fd^n=F\iota^{n+1}F\mathbf{d}^n$. Maka barisan-barisan eksak pendek tersebut
dapat dirakit kembali menjadi barisan eksak $(FX^\bullet,Fd^\bullet)$ di
$\mathcal{B}$.

(E2) $\implies$ (E3). Jika $F$ mempertahankan barisan eksak, ia memenuhi
kondisi (L2) dan (R2) dari Proposisi~\ref{prop:left-right-exact-functor}.

(E3) $\implies$ (E1). Misalkan
$0\to X'\to X\to X''\to0$ eksak di $\mathcal{A}$. Sekali lagi menerapkan
Proposisi~\ref{prop:left-right-exact-functor}, kita melihat bahwa
$0\to F(X')\to F(X)\to F(X'')$ dan
$F(X')\to F(X)\to F(X'')\to0$ keduanya eksak; jadi
$0\to F(X')\to F(X)\to F(X'')\to0$ eksak.
\end{bukti}

% segment-id: o014.li2.chapter2.exact-functors-injective-projective-objects.p004
Perhatikan bahwa (E3) $\iff$ (L3) $\wedge$ (R3).
% segment-id: o014.li2.chapter2.exact-functors-injective-projective-objects.d001
\begin{definisi}[Funktor eksak]\label{def:exact-functor}
	\index{eksak kiri, eksak kanan, eksak (left exact, right exact, exact)}
Untuk funktor aditif $F:\mathcal{A}\to\mathcal{B}$ antara kategori abelian,
pertimbangkan kondisi (L1)--(L3), (R1)--(R3) dari
Proposisi~\ref{prop:left-right-exact-functor}, serta kondisi (E1)--(E3) dari
Proposisi~\ref{prop:exact-functor}.
% segment-id: o014.li2.chapter2.exact-functors-injective-projective-objects.l005
\begin{itemize}
% segment-id: o014.li2.chapter2.exact-functors-injective-projective-objects.i012
	\item Jika salah satu dari (L1)--(L3) berlaku, maka $F$ disebut
	\emph{eksak kiri}.
% segment-id: o014.li2.chapter2.exact-functors-injective-projective-objects.i013
	\item Jika salah satu dari (R1)--(R3) berlaku, maka $F$ disebut
	\emph{eksak kanan}.
% segment-id: o014.li2.chapter2.exact-functors-injective-projective-objects.i014
	\item Jika salah satu dari (E1)--(E3) berlaku, maka $F$ disebut
	\emph{eksak}; ini ekuivalen dengan mengatakan bahwa $F$ eksak kiri sekaligus
	eksak kanan.
\end{itemize}
\end{definisi}

% segment-id: o014.li2.chapter2.exact-functors-injective-projective-objects.p005
Sebagai contoh, ekuivalensi antara kategori abelian tentu merupakan funktor eksak.
Contoh ekstrem lainnya adalah funktor nol: ia memetakan setiap objek ke objek nol
dan setiap morfisme ke morfisme nol; funktor ini juga eksak.

% segment-id: o014.li2.chapter2.exact-functors-injective-projective-objects.p006
Menurut Lema~\ref{prop:mono-epi-ker-coker}, funktor eksak kiri mempertahankan
monomorfisme dan funktor eksak kanan mempertahankan epimorfisme.

% segment-id: o014.li2.chapter2.exact-functors-injective-projective-objects.rem001
\begin{catatan}\label{rem:exactness-mono-epi}
Dengan mengingat (E1), jika $F$ diketahui eksak kiri (atau eksak kanan), maka
$F$ eksak jika dan hanya jika $F$ mempertahankan epimorfisme (atau
monomorfisme).
\end{catatan}

% segment-id: o014.li2.chapter2.exact-functors-injective-projective-objects.p007
Cara umum untuk memeriksa sifat eksak kiri/kanan adalah melalui funktor adjoin.

% segment-id: o014.li2.chapter2.exact-functors-injective-projective-objects.th001
\begin{teorema}\label{prop:exactness-adjoint}
Pertimbangkan sepasang funktor aditif antara kategori abelian $\mathcal{A}$ dan
$\mathcal{B}$:
% segment-id: o014.li2.chapter2.exact-functors-injective-projective-objects.g003
\[\begin{tikzcd}
	F: \mathcal{A} \arrow[r, yshift=0.3em] & \mathcal{B} :G \arrow[l, yshift=-0.3em] .
\end{tikzcd}\]
Misalkan $F$ dan $G$ dapat dilengkapi menjadi pasangan adjoin $(F,G,\varphi)$.
Maka $F$ eksak kanan dan $G$ eksak kiri; sesungguhnya, $F$ mempertahankan
$\varinjlim$ dan $G$ mempertahankan $\varprojlim$.
\end{teorema}

% segment-id: o014.li2.chapter2.exact-functors-injective-projective-objects.pf003
\begin{bukti}
Terapkan \cite[Teorema 2.8.12]{Li1} dan kondisi (L3), (R3) pada
Proposisi~\ref{prop:left-right-exact-functor}.
\end{bukti}

% segment-id: o014.li2.chapter2.exact-functors-injective-projective-objects.ex001
\begin{contoh}[Eksak kiri/kanan dari limit]\label{eg:limit-exactness}
Misalkan $\mathcal{A}$ merupakan kategori abelian, $I$ kategori sembarang, dan
anggap bahwa setiap funktor $\alpha:I\to\mathcal{A}$ memiliki $\varinjlim$ (atau
$\varprojlim$). Berdasarkan Proposisi~\ref{prop:functor-cat-Abel}, lengkapi
$\mathcal{A}^I$ dengan struktur kategori abelian. Kita akan menunjukkan bahwa
funktor $\varinjlim:\mathcal{A}^I\to\mathcal{A}$ (atau $\varprojlim$) eksak kanan
(atau eksak kiri).

Membahas kasus $\varinjlim$ saja sudah cukup. Berdasarkan Teorema~\ref{prop:exactness-adjoint},
salah satu strategi adalah menunjukkan bahwa $\varinjlim$ memiliki adjoin kanan.
Definisikan funktor diagonal $\Delta:\mathcal{A}\to\mathcal{A}^I$, yang memetakan
setiap $L\in\Obj(\mathcal{A})$ ke funktor konstan $\Obj(I)\ni i\mapsto L$. Kita
mempunyai pasangan adjoin
% segment-id: o014.li2.chapter2.exact-functors-injective-projective-objects.g004
\[\begin{tikzcd}
	\varinjlim : \mathcal{A}^I \arrow[shift left, r] & \mathcal{A}: \Delta \arrow[l, shift left] .
\end{tikzcd}\]
Ini merupakan konsekuensi langsung dari sifat universal $\varinjlim$: di
$\mathcal{A}^I$, memberikan morfisme dari funktor $\alpha:I\to\mathcal{A}$ ke
$\Delta(L)$ ekuivalen dengan memberikan keluarga morfisme kompatibel
$f_i:\alpha(i)\to L$; dengan kata lain, memberikan kerucut dengan dasar $\alpha$
dan puncak $L$. Lihat ulasan terkait pada \S\ref{sec:limit-functor}.
\end{contoh}

% segment-id: o014.li2.chapter2.exact-functors-injective-projective-objects.ex002
\begin{contoh}\label{eg:module-adjoint-pairs}
Misalkan $f:R\to S$ suatu homomorfisme gelanggang. Pertimbangkan funktor di
antara kategori modul kiri
% segment-id: o014.li2.chapter2.exact-functors-injective-projective-objects.g005
\[\begin{tikzcd}[column sep=large]
	R\dcate{Mod} \arrow[bend left=50, r, "{S \dotimes{R} (\cdot)}"] \arrow[bend right=50, r, "{\Hom_R({}_R S, \cdot)}"'] & S\dcate{Mod} \arrow[l, "{ {}_{R \to S} \mathcal{F} }"]
\end{tikzcd}\]
Di sini funktor pelupa ${}_{R\to S}\mathcal{F}$ tidak lain adalah pengubahan
suatu modul-$S$ menjadi modul-$R$ melalui $f$. Pembahasan funktor
$S\dotimes{R}(\cdot)$ dan $\Hom_R({}_R S,\cdot)$ dapat ditemukan di
\cite[\S 6.6]{Li1}; di sini ${}_R S$ berarti $S$ dipandang sebagai modul-$R$ kiri.
Menurut \cite[Korolari 6.6.8]{Li1},
% segment-id: o014.li2.chapter2.exact-functors-injective-projective-objects.q001
\[\left(S\dotimes{R}-,\;{}_{R\to S}\mathcal{F}\right)
\quad\text{dan}\quad
\left({}_{R\to S}\mathcal{F},\;\Hom_R({}_R S,-)\right)\]
merupakan pasangan adjoin. Teorema~\ref{prop:exactness-adjoint} kemudian menyiratkan
% segment-id: o014.li2.chapter2.exact-functors-injective-projective-objects.q002
\[S\dotimes{R}(\cdot)\ \text{eksak kanan},\quad
\Hom({}_R S,\cdot)\ \text{eksak kiri},\quad
{}_{R\to S}\mathcal{F}\ \text{eksak}.\]
Sifat eksak ini juga dapat diverifikasi langsung secara aljabar. Sebagai contoh,
untuk funktor pelupa ${}_{R\to S}\mathcal{F}$, apakah suatu barisan homomorfisme
modul $\cdots\to M^n\xrightarrow{d^n}M^{n+1}\to\cdots$ merupakan kompleks
(yaitu $d^{n+1}d^n=0$), atau eksak (yaitu
$\Image(d^n)=\Ker(d^{n+1})$), tidak bergantung pada perkalian di $R$ atau $S$,
melainkan hanya pada struktur grup aditif dari $M^n$. Dengan kata lain,
${}_{\Z\to S}\mathcal{F}:S\dcate{Mod}\to\cate{Ab}$ sudah merupakan funktor eksak.
Dari sudut pandang lain, kita dapat langsung memeriksa bahwa
${}_{R\to S}\mathcal{F}$ mempertahankan semua limit; hal ini dapat dicek langsung
dari konstruksi pada \cite[Teorema 6.2.2]{Li1}.
\end{contoh}
% segment-id: o014.li2.chapter2.exact-functors-injective-projective-objects.p008
Funktor eksak secara otomatis mempertahankan $\Image\simeq\Coim$ dari suatu
morfisme (Definisi~\ref{def:Im-Coim}, Proposisi~\ref{prop:Im-Coim-additive});
ia juga mempertahankan kohomologi.

% segment-id: o014.li2.chapter2.exact-functors-injective-projective-objects.pr003
\begin{proposisi}\label{prop:exact-preserves-cohomologies}
Misalkan $F:\mathcal{A}\to\mathcal{B}$ funktor eksak antara kategori abelian.
Untuk setiap kompleks $(X^\bullet,d^\bullet)$ di $\mathcal{A}$, terdapat
isomorfisme kanonik di $\mathcal{B}$
% segment-id: o014.li2.chapter2.exact-functors-injective-projective-objects.q003
\[F\Hm^n(X^\bullet,d^\bullet)\rightiso
\Hm^n(FX^\bullet,Fd^\bullet),\qquad n\in\Z.\]
\end{proposisi}

% segment-id: o014.li2.chapter2.exact-functors-injective-projective-objects.pf004
\begin{bukti}
Persoalan direduksi ke kasus kompleks tiga suku
$X'\xrightarrow{f}X\xrightarrow{g}X''$. Kembali ke definisi
\eqref{eqn:homology-3}:
% segment-id: o014.li2.chapter2.exact-functors-injective-projective-objects.q004
\begin{multline*}
	F\left(\Hm\left[X'\xrightarrow{f}X\xrightarrow{g}X''\right]\right)
	=F\Coker\left[\Image(f)\to\Ker(g)\right]\\
	\simeq\Coker\left[F\Image(f)\to F\Ker(g)\right]\\
	\simeq\Coker\left[\Image(Ff)\to\Ker(Fg)\right]
	=\Hm\left[FX'\xrightarrow{Ff}FX\xrightarrow{Fg}FX''\right],
\end{multline*}
dan semua isomorfisme yang muncul bersifat kanonik.
\end{bukti}

% segment-id: o014.li2.chapter2.exact-functors-injective-projective-objects.p009
Funktor eksak yang setia memiliki sifat-sifat yang sangat baik.

% segment-id: o014.li2.chapter2.exact-functors-injective-projective-objects.pr004
\begin{proposisi}\label{prop:faithful-exact}\index{funktor eksak setia (faithfully exact)}
Untuk funktor aditif $F:\mathcal{A}\to\mathcal{B}$ antara kategori abelian,
pernyataan berikut ekuivalen:
% segment-id: o014.li2.chapter2.exact-functors-injective-projective-objects.l006
\begin{enumerate}[(i)]
% segment-id: o014.li2.chapter2.exact-functors-injective-projective-objects.i015
	\item $F$ eksak dan setia;
% segment-id: o014.li2.chapter2.exact-functors-injective-projective-objects.i016
	\item $F$ eksak, dan untuk setiap $X\in\Obj(\mathcal{A})$ berlaku
	$FX=0\iff X=0$;
% segment-id: o014.li2.chapter2.exact-functors-injective-projective-objects.i017
	\item $X'\to X\to X''$ eksak di $\mathcal{A}$ jika dan hanya jika
	$FX'\to FX\to FX''$ eksak di $\mathcal{B}$.
\end{enumerate}
\end{proposisi}

% segment-id: o014.li2.chapter2.exact-functors-injective-projective-objects.pf005
\begin{bukti}
(i) $\implies$ (ii): Jika $FX=0$, maka
$F(\identity_X)=\identity_{FX}=0$, sehingga $\identity_X=0$ dan $X=0$.

(ii) $\implies$ (iii): Terapkan Proposisi~\ref{prop:exact-preserves-cohomologies}
pada kohomologi.

(iii) $\implies$ (i): Kondisi ini sudah mengimplikasikan bahwa $F$ eksak.
Misalkan $u:X\to Y$ memenuhi $Fu=0$. Karena
% segment-id: o014.li2.chapter2.exact-functors-injective-projective-objects.q005
\[X\xrightarrow{\identity}X\xrightarrow{u}Y\xrightarrow{\identity}Y,\]
citra barisan tersebut di bawah $F$ eksak, maka barisan itu sendiri eksak dan $u=0$.
\end{bukti}

% segment-id: o014.li2.chapter2.exact-functors-injective-projective-objects.p010
Bagi pembaca yang mengenal lokalisasi (lihat \S\ref{sec:cat-localization}), ini
memberikan contoh penting funktor eksak.

% segment-id: o014.li2.chapter2.exact-functors-injective-projective-objects.pr005
\begin{proposisi}[Keeksakan lokalisasi]\label{prop:Abel-cat-localization}
Misalkan $\mathcal{A}$ kategori abelian dan $S\subset\Mor(\mathcal{A})$ suatu
keluarga multiplikatif (Definisi~\ref{def:multiplicative-family}). Maka kategori
$\mathcal{A}[S^{-1}]$ yang diberikan oleh Teorema~\ref{prop:Gabriel-Zisman}
memiliki struktur kategori abelian kanonik, dan funktor lokalisasi
$Q:\mathcal{A}\to\mathcal{A}[S^{-1}]$ bersifat eksak.
\end{proposisi}

% segment-id: o014.li2.chapter2.exact-functors-injective-projective-objects.pf006
\begin{bukti}
Teorema~\ref{prop:localization-additivity} memberikan struktur kategori aditif
kanonik pada $\mathcal{A}[S^{-1}]$, sehingga $Q$ merupakan funktor aditif. Kita
akan menunjukkan bahwa setiap morfisme $f$ di $\mathcal{A}[S^{-1}]$ mempunyai
kernel dan kokernel serta merupakan morfisme ketat. Dari konstruksi
$\mathcal{A}[S^{-1}]$, setelah mengomposisikan $f$ dengan suatu isomorfisme dari
$S$, kita dapat mengasumsikan bahwa $f$ merupakan citra di bawah $Q$ dari suatu
morfisme di $\mathcal{A}$; hal ini tidak mengubah sifat yang hendak dibuktikan.
Lema~\ref{prop:localization-prod} menyatakan bahwa $Q$ mempertahankan semua
$\varinjlim$ dan $\varprojlim$ berhingga; khususnya, $Q$ memetakan
$\Ker$, $\Coker$, $\Image$, dan $\Coim$ di $\mathcal{A}$ ke konstruksi yang
bersesuaian di $\mathcal{A}[S^{-1}]$, sehingga juga mempertahankan diagram
\eqref{eqn:strict-morphism}. Ini menunjukkan bahwa $\mathcal{A}[S^{-1}]$
merupakan kategori abelian; kondisi (E3) pada Proposisi~\ref{prop:exact-functor}
menunjukkan bahwa $Q$ eksak.
\end{bukti}

% segment-id: o014.li2.chapter2.exact-functors-injective-projective-objects.p011
Jika lebih lanjut $\mathcal{A}$ merupakan kategori abelian $\Bbbk$-linear,
dengan $\Bbbk$ gelanggang komutatif, maka $Q$ merupakan funktor antara kategori
abelian $\Bbbk$-linear. Ini juga merupakan isi Teorema~\ref{prop:localization-additivity}.

% segment-id: o014.li2.chapter2.exact-functors-injective-projective-objects.p012
Contoh lain yang sangat penting adalah funktor $\Hom$. Jika $T$ objek dari
kategori abelian $\mathcal{A}$, maka $\Hom(T,\cdot)$ memberikan funktor
$\mathcal{A}\to\cate{Ab}$, sedangkan $\Hom(\cdot,T)$ memberikan funktor
$\mathcal{A}^{\opp}\to\cate{Ab}$; pada tingkat morfisme, keduanya memetakan
$f:X\to Y$ menjadi $f_*$ dan $f^*$ pada $\Hom$. Jelas keduanya merupakan funktor
aditif.

% segment-id: o014.li2.chapter2.exact-functors-injective-projective-objects.p013
Jika $\mathcal{A}$ kategori abelian $\Bbbk$-linear, funktor $\Hom$ dapat mengambil
nilai di $\Bbbk\dcate{Mod}$ dan menjadi funktor $\Bbbk$-linear; perluasan ini
sedikit memengaruhi pembahasan berikutnya sehingga tidak akan diuraikan tersendiri.

% segment-id: o014.li2.chapter2.exact-functors-injective-projective-objects.pr006
\begin{proposisi}[$\Hom$ eksak kiri]\label{prop:Hom-left-exact}
Misalkan $\mathcal{A}$ kategori abelian dan $T$ objek di $\mathcal{A}$. Maka
$\Hom(T,\cdot):\mathcal{A}\to\cate{Ab}$ dan
$\Hom(\cdot,T):\mathcal{A}^{\opp}\to\cate{Ab}$ keduanya merupakan funktor eksak kiri.
\end{proposisi}

% segment-id: o014.li2.chapter2.exact-functors-injective-projective-objects.pf007
\begin{bukti}
Dengan dualitas (menggantikan $\mathcal{A}$ dengan $\mathcal{A}^{\opp}$), cukup
menangani $\Hom(T,\cdot)$. Persoalannya adalah menunjukkan bahwa
$\Hom(T,\cdot)$ mempertahankan $\Ker$. Karena $\Ker$ di $\cate{Ab}$ hanyalah
kernel dalam teori grup, semuanya dapat diterjemahkan melalui sifat universal
$\Ker$.
\end{bukti}

% segment-id: o014.li2.chapter2.exact-functors-injective-projective-objects.d002
\begin{definisi}\label{def:injective-projective-obj}
	\index{objek injektif (injective)}
	\index{objek projektif (projective)}
Misalkan $X$ objek dari kategori abelian $\mathcal{A}$. Jika
$\Hom(X,\cdot):\mathcal{A}\to\cate{Ab}$ merupakan funktor eksak, maka $X$
disebut \emph{objek projektif}; jika
$\Hom(\cdot,X):\mathcal{A}^{\opp}\to\cate{Ab}$ merupakan funktor eksak, maka $X$
disebut \emph{objek injektif}.
\end{definisi}

% segment-id: o014.li2.chapter2.exact-functors-injective-projective-objects.p014
Menurut Catatan~\ref{rem:exactness-mono-epi} dan
Proposisi~\ref{prop:Hom-left-exact}, untuk menentukan apakah objek $X$ projektif
(atau injektif), cukup memeriksa apakah funktor $\Hom(X,\cdot)$ (atau
$\Hom(\cdot,X)$) mempertahankan epimorfisme. Karena itu:
% segment-id: o014.li2.chapter2.exact-functors-injective-projective-objects.l007
\begin{itemize}
% segment-id: o014.li2.chapter2.exact-functors-injective-projective-objects.i018
	\item Objek $P$ projektif jika dan hanya jika untuk setiap barisan eksak
	di $\mathcal{A}$ berbentuk $Y\xrightarrow{g}X\to0$ dan setiap morfisme
	$P\to X$, terdapat $P\to Y$ yang membuat diagram berikut komutatif.
% segment-id: o014.li2.chapter2.exact-functors-injective-projective-objects.g006
	\[\begin{tikzcd}
		& P \arrow[d] \arrow[ld, "\exists"'] & \\
	Y \arrow[r, "g"'] & X \arrow[r] & 0 \\
	\end{tikzcd}\]
	(ekuivalen dengan $g_*:\Hom(P,Y)\to\Hom(P,X)$ yang surjektif).
% segment-id: o014.li2.chapter2.exact-functors-injective-projective-objects.i019
	\item Objek $I$ injektif jika dan hanya jika untuk setiap barisan eksak
	di $\mathcal{A}$ berbentuk $0\to X\xrightarrow{f}Y$ dan setiap morfisme
	$X\to I$, terdapat $Y\to I$ yang membuat diagram berikut komutatif.
% segment-id: o014.li2.chapter2.exact-functors-injective-projective-objects.g007
	\[\begin{tikzcd}
		0 \arrow[r] & X \arrow[r, "f"] \arrow[d] & Y \arrow[ld, "\exists"] \\
		& I &
	\end{tikzcd}\]
	(ekuivalen dengan $f^*:\Hom(Y,I)\to\Hom(X,I)$ yang surjektif).
\end{itemize}

% segment-id: o014.li2.chapter2.exact-functors-injective-projective-objects.lm001
\begin{lema}\label{prop:inj-proj-split}
Pertimbangkan barisan eksak pendek
$0\to X'\xrightarrow{f}X\xrightarrow{g}X''\to0$ dalam kategori abelian.
Jika $X'$ objek injektif atau $X''$ objek projektif, barisan eksak pendek ini terbelah.
\end{lema}

% segment-id: o014.li2.chapter2.exact-functors-injective-projective-objects.pf008
\begin{bukti}
Dengan dualitas, tanpa mengurangi keumuman kita dapat mengasumsikan $X'$ objek
injektif. Dalam pembahasan di atas, pertimbangkan diagram komutatif
% segment-id: o014.li2.chapter2.exact-functors-injective-projective-objects.g008
\[\begin{tikzcd}
	0 \arrow[r] & X' \arrow[r, "f"] \arrow[d, "{\identity_{X'}}"'] & X \arrow[ld, "\exists r"] \\
	& X' &
\end{tikzcd}\]
kemudian substitusikan $rf=\identity_{X'}$ ke dalam Proposisi~\ref{prop:split-ses}.
\end{bukti}

% segment-id: o014.li2.chapter2.exact-functors-injective-projective-objects.lm002
\begin{lema}\label{prop:prod-injective-projective}
Pertimbangkan keluarga objek $(X_i)_{i\in I}$ dalam kategori abelian
$\mathcal{A}$. Jika koproduk $\coprod_{i\in I}X_i$ (atau produk
$\prod_{i\in I}X_i$) ada di $\mathcal{A}$, maka koproduk tersebut projektif
(atau produk tersebut injektif) jika dan hanya jika setiap $X_i$ demikian.
\end{lema}

% segment-id: o014.li2.chapter2.exact-functors-injective-projective-objects.pf009
\begin{bukti}
Dengan dualitas, cukup mempertimbangkan kasus koproduk
$\coprod_{i\in I}X_i$. Sifat universal memberikan isomorfisme funktor
% segment-id: o014.li2.chapter2.exact-functors-injective-projective-objects.q006
\[
\Hom_{\mathcal{A}}\left(\coprod_{i\in I}X_i,-\right)
\rightiso\prod_{i\in I}\Hom_{\mathcal{A}}\left(X_i,-\right):
\mathcal{A}\to\cate{Ab}.
\]
Semua persoalan kini mereduksi pada pengamatan elementer berikut: untuk keluarga
pemetaan $(f_i:A_i\to B_i)_{i\in I}$ dengan $A_i,B_i$ himpunan, pemetaan terinduksi
$(f_i)_{i\in I}:\prod_{i\in I}A_i\to\prod_{i\in I}B_i$ bersifat surjektif jika dan
hanya jika setiap $f_i$ surjektif.
\end{bukti}

% segment-id: o014.li2.chapter2.exact-functors-injective-projective-objects.p015
Sebagai contoh, ambil gelanggang $R$. Semua modul-$R$ bebas merupakan objek
projektif di $R\dcate{Mod}$. Memang, cukup menunjukkan bahwa $R$ sendiri
projektif; tetapi $\Hom(R,\cdot):R\dcate{Mod}\to\cate{Ab}$ isomorfik dengan funktor
pelupa $\mathcal{F}:R\dcate{Mod}\to\cate{Ab}$, yang memetakan homomorfisme
$\varphi:R\to X$ ke $\varphi(1)\in X$. Maka $\Hom(R,\cdot)$ eksak.

% segment-id: o014.li2.chapter2.exact-functors-injective-projective-objects.p016
Modul bebas sekaligus merupakan generator untuk $R\dcate{Mod}$; lihat Definisi
\ref{def:generators} dan Contoh~\ref{eg:generators}. Objek injektif yang juga
merupakan kogenerator (atau objek projektif yang juga generator) sangat berguna.
Berikut satu karakterisasi sederhana.

% segment-id: o014.li2.chapter2.exact-functors-injective-projective-objects.pr007
\begin{proposisi}[Kogenerator Injektif dan Generator Projektif]\label{prop:injective-cogenerator}
Dalam kategori abelian $\mathcal{A}$, objek injektif (atau projektif) $X$
merupakan kogenerator (atau generator) jika dan hanya jika untuk setiap
$T\in\Obj(\mathcal{A})$, $T\neq0$, berlaku $\Hom(T,X)\neq0$ (atau
$\Hom(X,T)\neq0$).

Hal ini ekuivalen pula dengan $\Hom(\cdot,X)$ (atau $\Hom(X,\cdot)$) yang
merupakan funktor eksak setia.
\end{proposisi}

% segment-id: o014.li2.chapter2.exact-functors-injective-projective-objects.pf010
\begin{bukti}
Kita hanya membahas kasus ketika $X$ injektif. Pertama, misalkan $X$ kogenerator.
Untuk
% segment-id: o014.li2.chapter2.exact-functors-injective-projective-objects.g009
\[\begin{tikzcd}
	T \arrow[r, shift left, "{\identity_T}"] \arrow[r, shift right, "0"'] & T
\end{tikzcd}\]
terapkan definisi kogenerator; maka terdapat $\delta\in\Hom(T,X)$ dengan
$\delta=\delta\circ\identity_T\neq\delta\circ0=0$.

Sekarang arah sebaliknya. Kita ingin menunjukkan bahwa jika $h:S\to T$ tak nol,
maka terdapat $\delta\in\Hom(T,X)$ sehingga $\delta h\neq0$. Perhatikan bahwa
terdapat $\delta'\in\Hom(\Image(h),X)$ dengan $\delta'\neq0$. Karena pertimbangan
mengenai epimorfisme, komposisi $S\twoheadrightarrow\Image(h)\xrightarrow{\delta'}X$
tak nol. Sekarang gunakan sifat injektif $X$ (lihat pembahasan setelah
Definisi~\ref{def:injective-projective-obj}) untuk memperluas $\delta'$ menjadi
$\delta:T\to X$.

Pernyataan tentang funktor eksak setia merupakan penerapan langsung
Proposisi~\ref{prop:faithful-exact}.
\end{bukti}

% segment-id: o014.li2.chapter2.exact-functors-injective-projective-objects.p017
Untuk menjalankan aljabar homologi pada kategori abelian, kita sering memerlukan
adanya cukup banyak objek injektif atau projektif.

% segment-id: o014.li2.chapter2.exact-functors-injective-projective-objects.d003
\begin{definisi}\label{def:enough-injectives}
	\index{cukup banyak objek injektif (enough injectives)}
	\index{cukup banyak objek projektif (enough projectives)}
Misalkan $\mathcal{A}$ kategori abelian.
% segment-id: o014.li2.chapter2.exact-functors-injective-projective-objects.l008
\begin{itemize}
% segment-id: o014.li2.chapter2.exact-functors-injective-projective-objects.i020
	\item Jika untuk setiap $X\in\Obj(\mathcal{A})$ terdapat objek injektif $I$ dan
	monomorfisme $X\hookrightarrow I$, maka $\mathcal{A}$ dikatakan memiliki
	\emph{cukup banyak objek injektif}.
% segment-id: o014.li2.chapter2.exact-functors-injective-projective-objects.i021
	\item Jika untuk setiap $X\in\Obj(\mathcal{A})$ terdapat objek projektif $P$ dan
	epimorfisme $P\twoheadrightarrow X$, maka $\mathcal{A}$ dikatakan memiliki
	\emph{cukup banyak objek projektif}.
\end{itemize}
\end{definisi}

% segment-id: o014.li2.chapter2.exact-functors-injective-projective-objects.p018
Kedua konsep ini saling dual. Cara umum untuk membangun objek injektif atau
projektif adalah menggunakan adjoin funktor eksak; berikut rinciannya.

% segment-id: o014.li2.chapter2.exact-functors-injective-projective-objects.pr008
\begin{proposisi}\label{prop:adjoint-injective-projective}
Pertimbangkan sepasang funktor antara kategori abelian
% segment-id: o014.li2.chapter2.exact-functors-injective-projective-objects.g010
\[\begin{tikzcd}
	\mathcal{A} \arrow[r, shift left, "F"] & \mathcal{B} \arrow[l, shift left, "G"]
\end{tikzcd}\]
dan asumsikan $F$ merupakan funktor eksak.
% segment-id: o014.li2.chapter2.exact-functors-injective-projective-objects.l009
\begin{itemize}
% segment-id: o014.li2.chapter2.exact-functors-injective-projective-objects.i022
	\item Jika $G$ adjoin kiri dari $F$, maka $G$ memetakan objek projektif di
	$\mathcal{B}$ ke objek projektif di $\mathcal{A}$;
% segment-id: o014.li2.chapter2.exact-functors-injective-projective-objects.i023
	\item jika $G$ adjoin kanan dari $F$, maka $G$ memetakan objek injektif di
	$\mathcal{B}$ ke objek injektif di $\mathcal{A}$.
\end{itemize}
\end{proposisi}

% segment-id: o014.li2.chapter2.exact-functors-injective-projective-objects.pf011
\begin{bukti}
Dengan dualitas, cukup mempertimbangkan kasus $G$ adjoin kiri. Dalam hal ini
$G$ pasti aditif (Korolari~\ref{prop:automatic-additivity}). Misalkan $P$ objek
projektif di $\mathcal{B}$. Untuk setiap barisan eksak $(X^\bullet,d^\bullet)$
di $\mathcal{A}$, terdapat isomorfisme kompleks di $\cate{Ab}$
% segment-id: o014.li2.chapter2.exact-functors-injective-projective-objects.q007
\[
\Hom_{\mathcal{A}}(GP,X^\bullet)\rightiso
\Hom_{\mathcal{B}}(P,FX^\bullet).
\]
Karena $F$ eksak, $(FX^\bullet,Fd^\bullet)$ merupakan barisan eksak sehingga
ruas kanan eksak di $\cate{Ab}$. Jadi
$\Hom_{\mathcal{A}}(GP,\cdot):\mathcal{A}\to\cate{Ab}$ merupakan funktor eksak.
\end{bukti}

% segment-id: o014.li2.chapter2.exact-functors-injective-projective-objects.p019
Sebagai contoh, pertimbangkan gelanggang $R$ dan funktor pelupa
$\mathcal{F}:R\dcate{Mod}\to\cate{Ab}$. Menurut Contoh~\ref{eg:module-adjoint-pairs},
$\mathcal{F}$ eksak dan memiliki adjoin kanan
$G:=\Hom_{\cate{Ab}}(R,\cdot)$. Proposisi~\ref{prop:adjoint-injective-projective}
menunjukkan bahwa $G$ memetakan objek injektif di $\cate{Ab}$ ke objek injektif
di $R\dcate{Mod}$; objek injektif di $\cate{Ab}$ mudah dikarakterisasi: objek
tersebut hanyalah modul-$\Z$ yang dapat dibagi. Ini merupakan cara standar dalam
teori modul untuk membangun modul-$R$ injektif dan menunjukkan bahwa
$R\dcate{Mod}$ memiliki cukup banyak objek injektif; lihat
\cite[Teorema 6.9.14]{Li1}.

% segment-id: o014.li2.chapter2.exact-functors-injective-projective-objects.p020
Di sisi lain, $R\dcate{Mod}$ juga memiliki cukup banyak objek projektif: untuk
setiap modul-$R$ $X$, ambil himpunan bagian $A\subset X$ yang menghasilkan $X$;
maka $R^{\oplus A}\twoheadrightarrow X$.

% segment-id: o014.li2.chapter2.exact-functors-injective-projective-objects.ex003
\begin{contoh}\label{eg:A2-injective-projective}
Latihan gabungan berikut melibatkan kategori abelian abstrak dan akan digunakan
di \S\sourcecrossref{sec:derived-primer}{3.12} untuk mempelajari barisan eksak panjang dari
funktor turunan. Misalkan $\mathcal{A}$ kategori abelian. Pertimbangkan kategori
$\mathbf{2}$ (diagramnya $0\to1$; lihat pengantar buku tentang kategori).
Kategori funktor $\mathcal{A}^{\mathbf{2}}$ tetap abelian
(Proposisi~\ref{prop:functor-cat-Abel}); ini adalah kategori panah: objeknya
merupakan morfisme $X_0\to X_1$ di $\mathcal{A}$, sedangkan morfismenya adalah
persegi komutatif di $\mathcal{A}$
% segment-id: o014.li2.chapter2.exact-functors-injective-projective-objects.g011
\[\begin{tikzcd}[row sep=small, column sep=small]
	X_0 \arrow[r] \arrow[d] & Y_0 \arrow[d] \\
	X_1 \arrow[r] & Y_1
\end{tikzcd}\]

% segment-id: o014.li2.chapter2.exact-functors-injective-projective-objects.p021
Untuk $i\in\{0,1\}$, funktor evaluasi
$\mathrm{ev}_i:\mathcal{A}^{\mathbf{2}}\to\mathcal{A}$ memetakan objek
$X_0\to X_1$ ke $X_i$. Selain itu, definisikan funktor aditif dari $\mathcal{A}$
ke $\mathcal{A}^{\mathbf{2}}$ melalui
\footnote{Catatan penerjemah (O014-C027): sumber menuliskan
$H:X\to[X\xrightarrow{\identity_X}X]$, sedangkan konteksnya mendefinisikan
pemetaan objek dari funktor $H$, sejajar dengan $L_0$ dan $L_1$. Target memakai
$\mapsto$ untuk memperbaiki panah penugasan tersebut.}
% segment-id: o014.li2.chapter2.exact-functors-injective-projective-objects.q008
\[L_0:X\mapsto[X\to0],\quad L_1:X\mapsto[0\to X],\quad
H:X\mapsto[X\xrightarrow{\identity_X}X],\]
untuk $X\in\Obj(\mathcal{A})$; definisi pada tingkat morfisme jelas. Kita akan
membuktikan hasil berikut.
% segment-id: o014.li2.chapter2.exact-functors-injective-projective-objects.l010
\begin{compactenum}[(i)]
% segment-id: o014.li2.chapter2.exact-functors-injective-projective-objects.i024
	\item Funktor $\mathrm{ev}_0$ dan $\mathrm{ev}_1$ keduanya eksak dan keduanya
	memetakan objek injektif (atau projektif) di $\mathcal{A}^{\mathbf{2}}$ ke objek
	injektif (atau projektif) di $\mathcal{A}$. Selain itu, $H$, $L_0$, dan $L_1$
	semuanya eksak.
% segment-id: o014.li2.chapter2.exact-functors-injective-projective-objects.i025
	\item Funktor $L_0,H$ (atau $H,L_1$) memetakan objek injektif (atau projektif)
	di $\mathcal{A}$ ke objek injektif (atau projektif) di $\mathcal{A}^{\mathbf{2}}$.
% segment-id: o014.li2.chapter2.exact-functors-injective-projective-objects.i026
	\item Jika $\mathcal{A}$ memiliki cukup banyak objek injektif (atau projektif),
	maka $\mathcal{A}^{\mathbf{2}}$ juga demikian.
\end{compactenum}

% segment-id: o014.li2.chapter2.exact-functors-injective-projective-objects.p022
Karena $\varinjlim$ dan $\varprojlim$ di $\mathcal{A}^{\mathbf{2}}$ dibangun suku
demi suku, funktor $\mathrm{ev}_0,\mathrm{ev}_1$ dan $H,L_0,L_1$ memang eksak.
Bagian yang tersisa dari (i) dan (ii) berasal dari hubungan adjoin yang dipenuhi
oleh $\mathrm{ev}_i$; diagramnya sebagai berikut, dan verifikasinya merupakan
latihan sederhana:
% segment-id: o014.li2.chapter2.exact-functors-injective-projective-objects.g012
\[\begin{tikzcd}[column sep=huge]
	\mathcal{A}^{\mathbf{2}} \arrow[r, "{\mathrm{ev}_0}" description] & \mathcal{A} \arrow[l, bend right, "{H: \text{adjoin kiri}}"'] \arrow[l, bend left, "L_0: \text{adjoin kanan}"] 
\end{tikzcd} \quad \begin{tikzcd}[column sep=huge]
	\mathcal{A}^{\mathbf{2}} \arrow[r, "{\mathrm{ev}_1}" description] & \mathcal{A} \arrow[l, bend right, "{L_1: \text{adjoin kiri}}"'] \arrow[l, bend left, "{H: \text{adjoin kanan}}"] .
\end{tikzcd}\]

% segment-id: o014.li2.chapter2.exact-functors-injective-projective-objects.p023
Untuk (iii), pertama-tama bahas kasus objek injektif. Diberikan objek
$[X_0\xrightarrow{f}X_1]$ di $\mathcal{A}^{\mathbf{2}}$, ambil monomorfisme
$\epsilon_i:X_i\hookrightarrow I_i$, dengan $I_i$ objek injektif,
$i\in\{0,1\}$. Dari sini diperoleh dua morfisme di $\mathcal{A}^{\mathbf{2}}$,
yang dapat ditulis sebagai diagram komutatif
% segment-id: o014.li2.chapter2.exact-functors-injective-projective-objects.g013
\[\begin{tikzcd}[column sep=large, /tikz/execute at end picture={
		\node[rectangle, draw, dashed, fit=(NE) (SE), inner xsep=0.5em, inner ysep=0em] {};
	}]
	X_0 \arrow[d, "f"'] \arrow[hookrightarrow, r, "\epsilon_0"] & |[alias=NE]| I_0 \arrow[d] \\
	X_1 \arrow[r] & |[alias=SE]| 0
\end{tikzcd} \quad \text{dan} \quad \begin{tikzcd}[column sep=large, /tikz/execute at end picture={
	\node[rectangle, draw, dashed, fit=(NE) (SE), inner xsep=0.5em, inner ysep=0em] {};
	}]
	X_0 \arrow[d, "f"'] \arrow[r, "\epsilon_1 f"] & |[alias=NE]| I_1 \arrow[d, "\identity"] \\
	X_1 \arrow[hookrightarrow, r, "\epsilon_1"'] & |[alias=SE]| I_1
\end{tikzcd}\]
Kedua kolom yang dibingkai masing-masing adalah $L_0(I_0)$ dan $H(I_1)$.
Menurut (ii), keduanya objek injektif di $\mathcal{A}^{\mathbf{2}}$, dan demikian pula
jumlah langsungnya (Lema~\ref{prop:prod-injective-projective}). Maka diagram
komutatif
% segment-id: o014.li2.chapter2.exact-functors-injective-projective-objects.g014
\[\begin{tikzcd}[column sep=large]
	X_0 \arrow[d, "f"'] \arrow[hookrightarrow, r, "{(\epsilon_0, \epsilon_1 f)}"] & I_0 \oplus I_1 \arrow[d, "\text{proyeksi}"] \\
	X_1 \arrow[hookrightarrow, r, "{\epsilon_1}"'] & I_1
\end{tikzcd}\]
membenamkan $[X_0\xrightarrow{f}X_1]$ ke dalam objek injektif. Untuk kasus objek
projektif, gunakan funktor $H$ dan $L_1$.

% segment-id: o014.li2.chapter2.exact-functors-injective-projective-objects.p024
\end{contoh}

% segment-id: o014.li2.chapter2.exact-functors-injective-projective-objects.p025
Pembahasan pada Contoh~\ref{eg:A2-injective-projective} dapat diperluas dari
$\mathcal{A}^{\mathbf{2}}$ ke kategori funktor umum $\mathcal{A}^{\mathcal{C}}$;
argumennya tidak mengandung kesulitan hakiki. Latihan bab ini akan menguraikannya.
